\subsection{DUAL OF A g-MODULE}

Let E, F be g-modules. Recall (Chap. I, §3, no. 3) that \(\operatorname{Hom}_{k}(\operatorname{E},\operatorname{F})\) has a canonical g-module structure. Let \(\varphi\) be an element of weight \(\lambda\) in \(\operatorname{Hom}_{k}(\operatorname{E},\operatorname{F})\) . If \(\mu \in h^{*}\) , then \(\varphi(\operatorname{E}^{\mu}) \subset \operatorname{F}^{\lambda + \mu}\) (Chap. VII, §1, no. 1, Prop. 2 (ii)). Thus, if E and F are finite dimensional, the elements of weight \(\lambda\) in \(\operatorname{Hom}_{k}(\operatorname{E},\operatorname{F})\) are the graded homomorphisms of degree \(\lambda\) in the sense of Algebra, Chap. II, §11, no. 2, Def. 4.

PROPOSITION 11. Let E be a finite dimensional g-module, and consider the g-module \(E^{*} = \operatorname{Hom}_{k}(E, k)\) .

\begin{enumerate}
\def\labelenumi{(\roman{enumi})}
\item
  An element \(\lambda \in \mathrm{P}\) is a weight of \(\mathrm{E}^*\) if and only if \(-\lambda\) is a weight of \(\mathrm{E}\) , and the multiplicity of \(\lambda\) in \(\mathrm{E}^*\) is equal to that of \(-\lambda\) in \(\mathrm{E}\) .
\item
  If \(\mathrm{E}\) is simple and has highest weight \(\omega\) , \(\mathrm{E}^*\) is simple and has highest weight \(-w_0(\omega)\) (cf.~no. 2, Remark 2).
\end{enumerate}

Consider k as a trivial g-module whose elements are of weight 0. By what was said above, the elements of \(E^{*}\) of weight \(\lambda\) are the homomorphisms from E to k which vanish on \(E^{\mu}\) if \(\mu \neq -\lambda\) . This proves (i). If E is simple, \(E^{*}\) is simple (Chap. I, §3, no. 3), and the last assertion follows from Remark 2 of no. 2.

Remarks. 1) Let E, E* be as in Prop. 11, and \(\sigma \in \mathrm{Aut}(\mathfrak{g},\mathfrak{h})\) be such that \(\varepsilon (\sigma) = -w_0\) in the notations of §5, no. 1 (§5, no. 2, Prop. 2). Let \(\rho ,\rho^{\prime}\) be the representations of \(\mathfrak{g}\) associated to E, E*. Then \(\rho \circ \sigma\) is a simple representation of \(\mathfrak{g}\) with highest weight \(-w_{0}(\omega)\) , so \(\rho \circ \sigma\) is equivalent to \(\rho^{\prime}\) .

\begin{enumerate}
\def\labelenumi{\arabic{enumi})}
\setcounter{enumi}{1}
\tightlist
\item
  Assume that \(w_{0} = -1\) . Then, for any finite dimensional g-module E, E is isomorphic to \(E^{*}\) . Recall that, if g is simple, \(w_{0} = -1\) in the following cases: g of type \(A_{1}, B_{l} (l \geq 2)\) , \(C_{l} (l \geq 2)\) , \(D_{l} (l \text{ even } \geq 4)\) , \(E_{7}, E_{8}, F_{4}, G_{2}\) (Chap. VI, Plates).
\end{enumerate}

Lemma 2. Let \(h^0 = \sum_{\alpha \in \mathrm{R}_{+}} H_\alpha\) . Then \(h^0 = \sum_{\alpha \in \mathrm{B}} a_\alpha H_\alpha\) , where the \(a_\alpha\) are integers \(\geq 1\) . Let \((b_\alpha)_{\alpha \in \mathrm{B}}, (c_\alpha)_{\alpha \in \mathrm{B}}\) be families of scalars such that \(b_\alpha c_\alpha = a_\alpha\) for all \(\alpha \in \mathrm{B}\) . Put \(x = \sum_{\alpha \in \mathrm{B}} b_\alpha X_\alpha, y = \sum_{\alpha \in \mathrm{B}} c_\alpha X_{-\alpha}\) . There exists a homomorphism \(\varphi\) from \(\mathfrak{sl}(2,k)\) to \(\mathfrak{g}\) such that \(\varphi(H) = h^0, \varphi(X_+) = x, \varphi(X_-) = y\) .

The fact that the \(a_{\alpha}\) are integers \(\geq 1\) follows from the fact that \((H_{\alpha})_{\alpha \in \mathrm{B}}\) is a basis of the root system \((H_{\alpha})_{\alpha \in \mathrm{B}}\) (cf.~Chap. VI, §1, no. 5, Remark 5). We have:

\[
\alpha (h ^ {0}) = 2\tag{1}
\]

for all \(\alpha \in \mathrm{B}\) (Chap. VI, §1, no. 10, Cor. of Prop. 29), so

\[
[ h ^ {0}, x ] = \sum_ {\alpha \in \mathrm{B}} b _ {\alpha} \alpha (h ^ {0}) X _ {\alpha} = 2 x\tag{2}
\]

\[
[ h ^ {0}, y ] = \sum_ {\alpha \in \mathrm{B}} c _ {\alpha} (- \alpha (h ^ {0})) X _ {- \alpha} = - 2 y.\tag{3}
\]

On the other hand,

\[
[ x, y ] = \sum_ {\alpha , \beta \in \mathrm{B}} b _ {\alpha} c _ {\beta} [ X _ {\alpha}, X _ {- \beta} ] = \sum_ {\alpha \in \mathrm{B}} b _ {\alpha} c _ {\alpha} [ X _ {\alpha}, X _ {- \alpha} ] = - \sum_ {\alpha \in \mathrm{B}} a _ {\alpha} H _ {\alpha} = - h ^ {0},\tag{4}
\]

hence the existence of the homomorphism \(\varphi\) .

PROPOSITION 12. Let E be a finite dimensional simple g-module, \(\omega\) its highest weight, and B the vector space of g-invariant bilinear forms on E. Let m be the integer \(\sum_{\alpha\in\mathrm{R}_{+}}\omega(H_{\alpha})\) , so that m/2 is the sum of the coordinates of \(\omega\) with respect to B (Chap. VI, §1, no. 10, Cor. of Prop. 29). Let \(w_{0}\) be the element of W such that \(w_{0}(B) = -B\) .

\begin{enumerate}
\def\labelenumi{(\roman{enumi})}
\item
  If \(w_0(\omega) \neq -\omega\) , then \(\mathcal{B} = 0\) .
\item
  Assume that \(w_{0}(\omega) = -\omega\) . Then B is of dimension 1, and every nonzero element of B is non-degenerate. If m is even (resp. odd), every element of B is symmetric (resp. alternating).
\end{enumerate}

\begin{enumerate}
\def\labelenumi{\alph{enumi})}
\item
  Let \(\Phi \in \mathcal{B}\) . The map \(\varphi\) from E to \(\mathrm{E}^*\) defined, for \(x, y \in \mathrm{E}\) , by \(\varphi(x)(y) = \Phi(x, y)\) is a homomorphism of \(\mathfrak{g}\) -modules. If \(\Phi \neq 0\) , then \(\varphi \neq 0\) , so \(\varphi\) is an isomorphism by Schur's lemma, and hence \(\Phi\) is non-degenerate. Consequently, the \(\mathfrak{g}\) -module E is isomorphic to the \(\mathfrak{g}\) -module \(\mathrm{E}^*\) , so that \(w_0(\omega) = -\omega\) . We have thus proved (i).
\item
  Assume from now on that \(w_{0}(\omega) = -\omega\) . Then E is isomorphic to \(E^{*}\) . The vector space B is isomorphic to \(\operatorname{Hom}_{\mathfrak{g}}(\mathrm{E}, \mathrm{E}^{*})\) , and hence to \(\operatorname{Hom}_{\mathfrak{g}}(\mathrm{E}, \mathrm{E})\) which is of dimension 1 (§6, no. 1, Prop. 1 (iii)). Hence \(\dim B = 1\) . Every non-zero element \(\Phi\) of B is non-degenerate by a). Put \(\Phi_{1}(x, y) = \Phi(y, x)\) for \(x, y \in E\) . By the preceding, there exists \(\lambda \in k\) such that \(\Phi_{1}(x, y) = \lambda \Phi(x, y)\) for all \(x, y \in E\) . Then \(\Phi(y, x) = \lambda \Phi(x, y) = \lambda^{2} \Phi(y, x)\) , so \(\lambda^{2} = 1\) and \(\lambda = \pm 1\) . Thus, \(\Phi\) is either symmetric or alternating.
\item
  By Chap. VII, §1, no. 3, Prop. 9 (v), \(E^{\lambda}\) and \(E^{\mu}\) are orthogonal with respect to \(\Phi\) if \(\lambda + \mu \neq 0\) . Since \(\Phi\) is non-degenerate, it follows that \(E^{\omega}, E^{-\omega}\) are not orthogonal with respect to \(\Phi\) .
\end{enumerate}

\(d\) ) There exists a homomorphism \(\varphi\) from \(\mathfrak{sl}(2,k)\) onto a subalgebra of \(\mathfrak{g}\) that takes \(H\) to \(\sum_{\alpha \in \mathrm{R}_{+}}H_{\alpha}\) (Lemma 2). Consider E as an \(\mathfrak{sl}(2,k)\) -module via this homomorphism. Then the elements of \(\mathrm{E}^{\lambda}\) are of weight \(\lambda \left(\sum_{\alpha \in \mathrm{R}_{+}}H_\alpha\right)\) . If \(\lambda \in \mathrm{P}\) is such that \(\mathrm{E}^{\lambda} \neq 0\) and \(\lambda \neq \omega, \lambda \neq -\omega\) , then \(-\omega < \lambda < \omega\) , so

\[
- m = - \omega \left(\sum_ {\alpha \in \mathrm{R} _ {+}} H _ {\alpha}\right) <   \lambda \left(\sum_ {\alpha \in \mathrm{R} _ {+}} H _ {\alpha}\right) <   \omega \left(\sum_ {\alpha \in \mathrm{R} _ {+}} H _ {\alpha}\right) = m.
\]

Let \(G\) be the isotypical component of type \(V(m)\) of the \(\mathfrak{sl}(2, k)\) -module \(E\) . By the preceding, \(G\) is of length 1 and contains \(E^{\omega}, E^{-\omega}\) . By \(c\) ), the restriction of \(\Phi\) to \(G\) is non-zero. By §1, no. 3, Remark 3, \(m\) is even or odd according as this restriction is symmetric or alternating. In view of \(b\) ), this completes the proof.

DEFINITION 2. A finite dimensional irreducible representation \(\rho\) of g is said to be orthogonal (resp. symplectic) if there exists on E a non-degenerate symmetric (resp. alternating) bilinear form invariant under \(\rho\) .

